\documentclass[10pt,a4paper]{article}
\usepackage{../../styles/neet_style}

\begin{document}

% Title Header
\begin{center}
  {\Huge \textbf{NEET UG 2027: Master Biology Chapter Notes}}\\[6pt]
  {\Large \textbf{Human Physiology: Respiration, Circulation, Excretion, Locomotion \& Coordination}}\\[4pt]
  {\small \textbf{Prepared, Compiled \& Authored by Satya Prakash} $\cdot$ Calibrated to Rationalized NCERT \& NMC Scope}
  \vspace{4pt}
  \hrule height 1.5pt
\end{center}

\pagestyle{fancy}
\fancyhf{}
\fancyhead[L]{\small \textbf{Human Physiology Master Notes}}
\fancyhead[R]{\small Satya Prakash}
\fancyfoot[L]{\small \textbf{NEET UG 2027 Prep Series}}
\fancyfoot[C]{\thepage}
\fancyfoot[R]{\small \textbf{Zoology Chapter Notes}}

\section{Breathing \& Exchange of Gases}

\begin{formulabox}[Respiratory Volumes \& Pulmonary Capacities Matrix]
All values represent healthy adult human averages measured by spirometer (residual volume cannot be measured by spirometer):
\begin{align*}
  \textbf{Tidal Volume (TV): } & 500\,\mathrm{mL} \quad (\approx 6000-8000\,\mathrm{mL/min} \text{ for healthy adult})\\
  \textbf{Inspiratory Reserve Volume (IRV): } & 2500 - 3000\,\mathrm{mL}\\
  \textbf{Expiratory Reserve Volume (ERV): } & 1000 - 1100\,\mathrm{mL}\\
  \textbf{Residual Volume (RV): } & 1100 - 1200\,\mathrm{mL} \quad (\text{air remaining after maximal expiration})\\
  \textbf{Inspiratory Capacity (IC): } & \mathrm{TV} + \mathrm{IRV} = 3000 - 3500\,\mathrm{mL}\\
  \textbf{Expiratory Capacity (EC): } & \mathrm{TV} + \mathrm{ERV} = 1500 - 1600\,\mathrm{mL}\\
  \textbf{Functional Residual Capacity (FRC): } & \mathrm{ERV} + \mathrm{RV} = 2100 - 2300\,\mathrm{mL}\\
  \textbf{Vital Capacity (VC): } & \mathrm{ERV} + \mathrm{TV} + \mathrm{IRV} = 3500 - 4500\,\mathrm{mL}\\
  \textbf{Total Lung Capacity (TLC): } & \mathrm{RV} + \mathrm{ERV} + \mathrm{TV} + \mathrm{IRV} = \mathrm{VC} + \mathrm{RV} = 5000 - 6000\,\mathrm{mL}
\end{align*}
\end{formulabox}

\begin{conceptbox}[Partial Pressures \& Gas Exchange Gradients]
Diffusion of gases across the three-layered respiratory membrane (alveolar squamous epithelium, basement membrane, capillary endothelium) occurs along partial pressure gradients:
\begin{center}
\small
\begin{tabular}{|l|c|c|c|c|c|}
\hline
\textbf{Gas} & \textbf{Atmospheric Air} & \textbf{Alveoli} & \textbf{Deoxygenated Blood} & \textbf{Oxygenated Blood} & \textbf{Tissues} \\
\hline
$\mathrm{O_2}$ & $159\,\mathrm{mmHg}$ & $\mathbf{104\,\mathrm{mmHg}}$ & $40\,\mathrm{mmHg}$ & $\mathbf{95\,\mathrm{mmHg}}$ & $40\,\mathrm{mmHg}$ \\
$\mathrm{CO_2}$ & $0.3\,\mathrm{mmHg}$ & $\mathbf{40\,\mathrm{mmHg}}$ & $45\,\mathrm{mmHg}$ & $\mathbf{40\,\mathrm{mmHg}}$ & $45\,\mathrm{mmHg}$ \\
\hline
\end{tabular}
\end{center}
\textbf{Solubility Factor:} The solubility of $\mathrm{CO}_2$ is \textbf{20--25 times higher} than that of $\mathrm{O}_2$. Hence, the amount of $\mathrm{CO}_2$ that can diffuse across the diffusion membrane per unit difference in partial pressure is much higher. Total membrane thickness is $< 1\,\mathrm{mm}$ (around $0.2\,\mathrm{mm}$).
\end{conceptbox}

\begin{conceptbox}[Oxygen Dissociation Curve \& Gas Transport Mechanisms]
\textbf{1. Transport of Oxygen:}
\begin{itemize}
  \item $97\%$ bound to hemoglobin as \textbf{Oxyhemoglobin} in RBCs ($1\,\mathrm{g}$ Hb binds $1.34\,\mathrm{mL}\; \mathrm{O}_2$; each Hb molecule binds max $4$ molecules of $\mathrm{O}_2$).
  \item $3\%$ dissolved in blood plasma.
  \item \textbf{Delivery Yield:} Under normal physiological conditions, every $100\,\mathrm{mL}$ of oxygenated blood delivers approximately \textbf{$5\,\mathrm{mL}$ of $\mathrm{O}_2$} to the tissues.
\end{itemize}
\textbf{2. Oxygen-Hemoglobin Dissociation Curve (Sigmoid / S-shaped):}
\begin{itemize}
  \item \textbf{Shift to Right (Bohr Effect / $\mathrm{O}_2$ Release at Tissues):} Promoted by:
  \begin{equation}
    \uparrow p\mathrm{CO_2}, \quad \uparrow [\mathrm{H}^+] \ (\downarrow \mathrm{pH}), \quad \uparrow \text{Temperature}, \quad \uparrow 2,3\text{-BPG}
  \end{equation}
  \item \textbf{Shift to Left (Loading at Alveoli):} Promoted by $\uparrow p\mathrm{O_2}, \downarrow p\mathrm{CO_2}, \downarrow [\mathrm{H}^+] \ (\uparrow \mathrm{pH}), \downarrow \text{Temperature}$.
\end{itemize}
\textbf{3. Transport of Carbon Dioxide:}
\begin{itemize}
  \item As \textbf{Bicarbonate ions ($\mathrm{HCO}_3^-$):} $\mathbf{70\%}$ facilitated by RBC enzyme \textit{carbonic anhydrase} and chloride shift (Hamburger phenomenon).
  \item As \textbf{Carbamino-hemoglobin:} $\mathbf{20-25\%}$ bound to amino groups of hemoglobin.
  \item In \textbf{Dissolved state} in plasma: $\mathbf{7\%}$.
  \item \textbf{Delivery Yield:} Every $100\,\mathrm{mL}$ of deoxygenated blood delivers approximately \textbf{$4\,\mathrm{mL}$ of $\mathrm{CO}_2$} to the alveoli.
\end{itemize}
\end{conceptbox}

\begin{trapbox}[NEET Trap: Respiratory Regulation \& Clinical Pathology]
\begin{itemize}
  \item \textbf{Primary Chemical Sensor:} Chemosensitive area located adjacent to the rhythm center in medulla is \textbf{highly sensitive to $\mathrm{CO}_2$ and $\mathrm{H}^+$ ions}. The role of oxygen in the regulation of respiratory rhythm is \textbf{quite insignificant}!
  \item \textbf{Respiratory Centers:} Respiratory rhythm center is in \textbf{Medulla oblongata}. \textbf{Pneumotaxic center} is in \textbf{Pons varolii} (moderates rhythm center and reduces duration of inspiration).
  \item \textbf{Asthma:} Allergic inflammation of bronchi and bronchioles causing wheezing.
  \item \textbf{Emphysema:} Chronic disorder in which \textbf{alveolar walls are damaged}, decreasing respiratory surface area; majorly caused by cigarette smoking ($\alpha$-1-antitrypsin used in treatment).
  \item \textbf{Occupational Respiratory Disorders:} Silicosis, asbestosis; chronic exposure leads to \textbf{fibrosis} (proliferation of fibrous tissues) in upper parts of lungs.
\end{itemize}
\end{trapbox}

\section{Body Fluids \& Circulation}

\begin{conceptbox}[Blood Composition, Formed Elements \& Coagulation Cascade]
Blood is a specialized fluid connective tissue consisting of plasma ($55\%$) and formed elements ($45\%$):
\begin{itemize}
  \item \textbf{Plasma:} $90-92\%$ water, $6-8\%$ proteins (Fibrinogen for clotting, Globulins for defense/immunity, Albumins for osmotic balance).
  \item \textbf{Erythrocytes (RBCs):} Biconcave, enucleated in adult mammals, $5.0-5.5 \times 10^6/\mathrm{mm}^3$. Average lifespan = $120$ days; destroyed in the \textbf{spleen} (graveyard of RBCs). Hemoglobin count = $12-16\,\mathrm{g}/100\,\mathrm{mL}$.
  \item \textbf{Leukocytes (WBCs):} Nucleated, $6000-8000/\mathrm{mm}^3$.
  \begin{itemize}
    \item \textit{Granulocytes:} \textbf{Neutrophils} ($60-65\%$, phagocytic), \textbf{Eosinophils} ($2-3\%$, resist infections, allergic response), \textbf{Basophils} ($0.5-1\%$, secrete histamine, serotonin, heparin; mediate inflammation).
    \item \textit{Agranulocytes:} \textbf{Monocytes} ($6-8\%$, phagocytic), \textbf{Lymphocytes} ($20-25\%$, B and T lymphocytes for adaptive immunity).
  \end{itemize}
  \item \textbf{Thrombocytes (Platelets):} Cell fragments produced from \textbf{megakaryocytes} in bone marrow ($1.5-3.5 \times 10^5/\mathrm{mm}^3$).
  \item \textbf{Blood Coagulation Cascade (Enzyme Cascade Theory):}
  \begin{equation}
    \text{Thromboplastin (Injured Tissue + Platelets)} \xrightarrow{\mathrm{Ca}^{2+}} \text{Prothrombinase / Thrombokinase}
  \end{equation}
  \begin{equation}
    \text{Prothrombin (Inactive)} \xrightarrow[\mathrm{Ca}^{2+}]{\text{Thrombokinase}} \text{Thrombin (Active)}
  \end{equation}
  \begin{equation}
    \text{Fibrinogen (Soluble)} \xrightarrow{\text{Thrombin}} \text{Fibrin (Insoluble Mesh)} \xrightarrow{\text{traps dead elements}} \text{Blood Clot}
  \end{equation}
  \textbf{Calcium ($\mathrm{Ca}^{2+}$)} plays an indispensable cofactor role in almost all clotting steps!
\end{itemize}
\end{conceptbox}

\begin{formulabox}[Cardiac Cycle, Hemodynamics \& Electrocardiogram (ECG)]
\textbf{1. Cardiac Cycle Timing ($72\,\text{beats/min} \implies \text{Duration} = 0.8\,\mathrm{sec}$):}
\begin{align*}
  \text{Atrial Systole: } & 0.1\,\mathrm{sec} \quad (\text{increases ventricular filling by } 30\%)\\
  \text{Ventricular Systole: } & 0.3\,\mathrm{sec} \quad (\text{isovolumetric contraction } \to \text{ ejection period})\\
  \text{Joint Diastole: } & 0.4\,\mathrm{sec} \quad (\text{passive filling of atria and ventricles})\\
  \textbf{Stroke Volume (SV): } & \mathrm{EDV} - \mathrm{ESV} = 120\,\mathrm{mL} - 50\,\mathrm{mL} = \mathbf{70\,\mathrm{mL}}\\
  \textbf{Cardiac Output (CO): } & \text{Stroke Volume} \times \text{Heart Rate} = 70\,\mathrm{mL} \times 72/\mathrm{min} \approx \mathbf{5040\,\mathrm{mL/min} \approx 5\,\mathrm{L/min}}
\end{align*}
\textbf{2. Heart Sounds:}
\begin{itemize}
  \item \textbf{First Heart Sound (LUBB):} Closure of \textbf{atrioventricular valves} (Tricuspid and Bicuspid) at the onset of ventricular systole; low pitched, longer duration.
  \item \textbf{Second Heart Sound (DUBB):} Closure of \textbf{semilunar valves} at the onset of ventricular diastole; high pitched, shorter duration.
\end{itemize}
\textbf{3. Standard Electrocardiogram (ECG) Deflections:}
\begin{itemize}
  \item \textbf{P-wave:} Depolarization of atria (leads to atrial contraction).
  \item \textbf{QRS Complex:} Depolarization of ventricles (marks onset of ventricular contraction; counting QRS complexes per unit time gives heart rate).
  \item \textbf{T-wave:} Repolarization of ventricles (return to resting state; end of T-wave marks end of systole).
\end{itemize}
\end{formulabox}

\begin{conceptbox}[Double Circulation \& Hepatic Portal System]
\begin{itemize}
  \item \textbf{Pulmonary Circulation:} Deoxygenated blood from Right Ventricle $\xrightarrow{\text{Pulmonary Artery}}$ Lungs $\xrightarrow{\text{Pulmonary Veins}}$ Left Atrium.
  \item \textbf{Systemic Circulation:} Oxygenated blood from Left Ventricle $\xrightarrow{\text{Aorta}}$ Body Tissues $\xrightarrow{\text{Vena Cava}}$ Right Atrium.
  \item \textbf{Hepatic Portal System:} Unique vascular connection carrying blood from the \textbf{digestive tract to the liver} via the Hepatic Portal Vein before reaching systemic circulation.
  \item \textbf{Coronary Circulation:} Exclusive network of blood vessels supplying muscular walls of the heart.
\end{itemize}
\end{conceptbox}

\section{Excretory Products \& Their Elimination}

\begin{conceptbox}[Modes of Nitrogenous Excretion \& Nephron Anatomy]
\begin{itemize}
  \item \textbf{Ammonotelism:} Excretion of toxic ammonia requiring large volumes of water; diffuses directly across body surfaces or gill surfaces as ammonium ions (e.g., bony fishes, aquatic amphibians, aquatic insects). Kidneys play no significant role in ammonia removal.
  \item \textbf{Ureotelism:} Urea synthesized in the liver via Ornithine cycle; requires moderate water (e.g., mammals, terrestrial amphibians, marine cartilaginous fishes).
  \item \textbf{Uricotelism:} Uric acid excreted as pellet/paste with minimal water loss; least toxic (e.g., reptiles, birds, land snails, insects).
  \item \textbf{Nephron Types:}
  \begin{itemize}
    \item \textbf{Cortical Nephrons ($85\%$):} Henle's loop is short and extends only very little into medulla; vasa recta absent or highly reduced.
    \item \textbf{Juxtamedullary Nephrons ($15\%$):} Henle's loop is extremely long and runs deep into the inner medulla; well-developed vasa recta; \textbf{essential for concentrating urine}.
  \end{itemize}
\end{itemize}
\end{conceptbox}

\begin{formulabox}[Renal Physiology: GFR \& Counter-Current Mechanism]
\textbf{1. Glomerular Filtration Rate (GFR):}
\begin{equation}
  \text{Renal Blood Flow} \approx 1100-1200\,\mathrm{mL/min} \quad (\approx 20-25\% \text{ of Cardiac Output})
\end{equation}
\begin{equation}
  \mathbf{GFR = 125\,\mathrm{mL/min} = 180\,\mathrm{litres/day}} \implies \text{Urine output } \approx \mathbf{1.5\,\mathrm{litres/day}} \ (99\% \text{ reabsorbed!})
\end{equation}
\textbf{2. Net Filtration Pressure (NFP):}
\begin{equation}
  \mathrm{NFP = GBHP - (BCOP + CHP)} = 60 - (32 + 18) = \mathbf{10\,\mathrm{mmHg}}
\end{equation}
\textbf{3. Counter-Current Multiplier \& Exchanger:}
\begin{itemize}
  \item Operated between \textbf{Loop of Henle} and \textbf{Vasa Recta}.
  \item Flow of filtrate in two limbs of Henle's loop is in opposite directions; flow of blood in two limbs of vasa recta is also in opposite directions.
  \item \textbf{Descending Limb of Henle:} Permeable to water, almost impermeable to electrolytes $\implies$ filtrate becomes concentrated as it moves down ($300 \to 1200\,\mathrm{mOsmol/L}$).
  \item \textbf{Ascending Limb of Henle:} Impermeable to water, actively/passively transports electrolytes ($\mathrm{NaCl}$) into medullary interstitium $\implies$ filtrate becomes diluted as it ascends ($1200 \to 200\,\mathrm{mOsmol/L}$).
  \item Creates an osmolarity gradient in medullary interstitium: \textbf{$300\,\mathrm{mOsmol/L}$ in cortex $\to 1200\,\mathrm{mOsmol/L}$ in inner medulla}.
  \item Human kidneys can produce urine nearly \textbf{4 times more concentrated} than the initial filtrate ($1200\,\mathrm{mOsmol/L}$ vs $300\,\mathrm{mOsmol/L}$).
\end{itemize}
\end{formulabox}

\begin{conceptbox}[Hormonal Regulation of Kidney Function]
\begin{enumerate}[leftmargin=*]
  \item \textbf{RAAS (Renin-Angiotensin-Aldosterone System):}
  \begin{equation}
    \downarrow \text{GFR / Blood Pressure} \xrightarrow{\text{detected by JGA}} \text{Juxtaglomerular cells release } \mathbf{Renin}
  \end{equation}
  \begin{equation}
    \text{Angiotensinogen (Liver)} \xrightarrow{\text{Renin}} \text{Angiotensin I} \xrightarrow{\text{ACE}} \text{Angiotensin II (Potent Vasoconstrictor)}
  \end{equation}
  Angiotensin II increases GFR directly by arteriolar constriction AND stimulates \textbf{adrenal cortex} to secrete \textbf{Aldosterone} $\implies$ promotes active reabsorption of $\mathrm{Na}^+$ and water from DCT $\implies \uparrow$ Blood Volume and BP.
  \item \textbf{Antidiuretic Hormone (ADH / Vasopressin):} Released by posterior pituitary when hypothalamic osmoreceptors detect $\uparrow$ blood osmolarity (dehydration); makes DCT and Collecting Duct permeable to water via aquaporins, preventing diuresis.
  \item \textbf{Atrial Natriuretic Factor (ANF):} Released by atrial walls of heart in response to $\uparrow$ blood flow/pressure; causes \textbf{vasodilation} and promotes natriuresis ($\mathrm{Na}^+$ excretion). Acts as a negative feedback check on RAAS!
\end{enumerate}
\end{conceptbox}

\section{Locomotion \& Movement: Musculoskeletal Mechanics}

\begin{conceptbox}[Sarcomere Architecture \& Sliding Filament Theory]
A sarcomere is the basic functional contractile unit of a muscle fiber, bounded between \textbf{two successive Z-lines}:
\begin{itemize}
  \item \textbf{A-Band (Anisotropic):} Contains both thick (myosin) and thin (actin) filaments; length \textbf{remains completely constant} during contraction.
  \item \textbf{I-Band (Isotropic):} Contains only thin actin filaments; bisected by elastic \textbf{Z-line}; \textbf{shortens} during contraction.
  \item \textbf{H-Zone:} Central region of A-band where thick filaments are not overlapped by thin filaments; bisected by fibrous \textbf{M-line}; \textbf{disappears} during maximal contraction.
\end{itemize}
\textbf{Myofilament Biochemistry:}
\begin{itemize}
  \item \textbf{Thin Filament:} Composed of 2 helically wound polymeric \textbf{F-actin} filaments, 2 filamentous \textbf{Tropomyosin} running along F-actin grooves, and complex \textbf{Troponin} triplets (TnT, TnI, TnC). \textbf{Troponin-C (TnC)} binds $\mathrm{Ca}^{2+}$, unmasking active myosin-binding sites on actin.
  \item \textbf{Thick Filament:} Composed of polymerized \textbf{Myosin} monomer units (\textbf{Meromyosins}):
  \begin{itemize}
    \item \textbf{Heavy Meromyosin (HMM):} Globular head and short arm; projects outwards as cross-bridge; contains actin-binding site and active \textbf{ATPase enzyme site}.
    \item \textbf{Light Meromyosin (LMM):} Tail portion.
  \end{itemize}
\end{itemize}
\textbf{Contraction Cycle:} Action potential depolarizes sarcolemma $\to$ Sarcoplasmic reticulum releases $\mathrm{Ca}^{2+} \to \mathrm{Ca}^{2+}$ binds to Troponin-C $\to$ Tropomyosin shifts $\to$ Myosin head hydrolyzes ATP ($\mathrm{ADP} + \mathrm{P_i}$) and binds actin forming cross-bridge $\to$ \textbf{Power Stroke} pulls thin filaments towards center of A-band $\to$ New ATP binds myosin head, causing detachment.
\end{conceptbox}

\begin{conceptbox}[Human Skeletal Framework (206 Bones) \& Articulations]
\begin{itemize}
  \item \textbf{Axial Skeleton (80 Bones):}
  \begin{itemize}
    \item \textbf{Skull (22 bones):} Cranial (8) $+$ Facial (14). Associated: Hyoid (1) $+$ Ear Ossicles ($2 \times 3 = 6$: Malleus, Incus, Stapes).
    \item \textbf{Vertebral Column (26 in adult):} Cervical (7, $C_1=\text{Atlas}, C_2=\text{Axis}$), Thoracic (12), Lumbar (5), Sacral (1 fused), Coccygeal (1 fused).
    \item \textbf{Thoracic Cage:} Sternum (1) $+$ Ribs (12 pairs = 24). True ribs ($1^{\mathrm{st}}-7^{\mathrm{th}}$ pair, vertebrosternal), False ribs ($8^{\mathrm{th}}, 9^{\mathrm{th}}, 10^{\mathrm{th}}$ pair, vertebrochondral), Floating ribs ($11^{\mathrm{th}}, 12^{\mathrm{th}}$ pair, vertebral).
  \end{itemize}
  \item \textbf{Appendicular Skeleton (126 Bones):}
  \begin{itemize}
    \item \textbf{Forelimb ($30 \times 2 = 60$):} Humerus (1), Radius (1), Ulna (1), Carpals (8 wrist), Metacarpals (5 palm), Phalanges (14 digits, formula 2,3,3,3,3).
    \item \textbf{Pectoral Girdle ($2 \times 2 = 4$):} Clavicle (collar bone) $+$ Scapula (has acromion process, spine, and glenoid cavity for head of humerus).
    \item \textbf{Hindlimb ($30 \times 2 = 60$):} Femur (thigh, longest/strongest), Patella (kneecap, sesamoid bone), Tibia (1), Fibula (1), Tarsals (7 ankle), Metatarsals (5), Phalanges (14).
    \item \textbf{Pelvic Girdle ($2$ coxal bones):} Each formed by fusion of Ilium, Ischium, and Pubis; meet ventrally at pubic symphysis (fibrous cartilage); cavity = \textbf{Acetabulum} (articulates with femur).
  \end{itemize}
  \item \textbf{Joint Classifications:}
  \begin{itemize}
    \item \textbf{Fibrous (Synarthroses):} Immovable (sutures between flat skull bones).
    \item \textbf{Cartilaginous (Amphiarthroses):} Limited movement (pubic symphysis, between adjacent vertebrae).
    \item \textbf{Synovial (Diarthroses):} Fluid-filled synovial cavity; freely movable.
    \begin{itemize}
      \item \textit{Ball and Socket:} Shoulder (glenoid-humerus) and Hip (acetabulum-femur).
      \item \textit{Hinge Joint:} Knee joint, Elbow joint, Interphalangeal joints.
      \item \textit{Pivot Joint:} Between Atlas and Axis ($C_1-C_2$).
      \item \textit{Gliding Joint:} Between adjacent Carpals.
      \item \textit{Saddle Joint:} Between Carpal and Metacarpal of human thumb.
    \end{itemize}
  \end{itemize}
\end{itemize}
\end{conceptbox}

\section{Neural Control \& Chemical Coordination}

\begin{formulabox}[Nerve Impulse: Resting, Depolarization \& Repolarization]
\begin{align*}
  \textbf{Resting Membrane Potential: } & \mathbf{-70\,\mathrm{mV}} \quad (\text{Inside negative; permeable to } \mathrm{K}^+, \text{ impermeable to } \mathrm{Na}^+)\\
  \textbf{Sodium-Potassium Pump: } & \text{Transports } \mathbf{3\,\mathrm{Na}^+} \text{ out for } \mathbf{2\,\mathrm{K}^+} \text{ in} \ (\text{Consumes 1 ATP})\\
  \textbf{Action Potential: } & \mathbf{+30\,\mathrm{mV}} \quad (\text{Voltage-gated } \mathrm{Na}^+ \text{ channels open} \implies \mathrm{Na}^+ \text{ influx})\\
  \textbf{Repolarization: } & \mathrm{Na}^+ \text{ channels close, voltage-gated } \mathrm{K}^+ \text{ channels open} \implies \mathrm{K}^+ \text{ efflux}\\
  \textbf{Synapse: } & \text{AP at terminal} \to \mathrm{Ca}^{2+} \text{ influx} \to \text{ACh exocytosis} \to \text{receptors}
\end{align*}
\end{formulabox}

\begin{conceptbox}[Endocrine Master Table: Glands, Hormones, Target \& Disorders]
\begin{center}
\scriptsize
\begin{tabular}{|p{2.2cm}|p{2.5cm}|p{1.8cm}|p{7.0cm}|}
\hline
\textbf{Endocrine Gland} & \textbf{Hormone} & \textbf{Nature} & \textbf{Major Action \& Pathology} \\
\hline
Anterior Pituitary & Growth Hormone (GH) & Peptide & Childhood excess: \textbf{Gigantism}; adult excess: \textbf{Acromegaly}; deficiency: \textbf{Pituitary Dwarfism}. \\
Posterior Pituitary & Oxytocin, Vasopressin & Peptide & Oxytocin: uterine contraction \& milk ejection. Vasopressin: water reabsorption; deficiency $\implies$ \textbf{Diabetes Insipidus}. \\
Thyroid & Thyroxine ($T_4, T_3$) & Iodinated & BMR regulation. Pregnancy deficiency $\implies$ \textbf{Cretinism}; adult: \textbf{Myxoedema}; excess: \textbf{Graves' disease}. \\
Thyroid / C-cells & Thyrocalcitonin (TCT) & Peptide & \textbf{Hypocalcemic}; lowers blood $\mathrm{Ca}^{2+}$ via bone deposition. \\
Parathyroid & Parathormone (PTH) & Peptide & \textbf{Hypercalcemic}; bone resorption; antagonist to TCT. \\
Adrenal Cortex & Aldosterone, Cortisol & Steroid & $\mathrm{Na}^+$ retention. Hyposecretion $\implies$ \textbf{Addison's disease}; hypersecretion $\implies$ \textbf{Cushing's syndrome}. \\
Adrenal Medulla & Adrenaline & Catecholamine & Emergency / Fight-or-Flight; $\uparrow$ heart rate, glycogenolysis. \\
Pancreas ($\beta$-cells) & Insulin & Peptide & Hypoglycemic; promotes cellular glucose uptake; deficiency $\implies$ \textbf{Diabetes Mellitus}. \\
Pancreas ($\alpha$-cells) & Glucagon & Peptide & Hyperglycemic; stimulates glycogenolysis; antagonist to insulin. \\
\hline
\end{tabular}
\end{center}
\end{conceptbox}

\begin{ncertbox}[Mechanism of Hormone Action: Peptide vs Steroid Pathway]
\begin{itemize}
  \item \textbf{Peptide / Protein / Amine Hormones (Water-Soluble):} Insulin, glucagon, pituitary hormones, epinephrine. Cannot cross lipid bilayer; bind to \textbf{extracellular membrane-bound receptors} $\implies$ generate second messengers (\textbf{cyclic AMP, $\mathrm{IP}_3, \mathrm{Ca}^{2+}$}) $\implies$ cascade modifies cellular biochemical metabolism.
  \item \textbf{Steroid \& Iodothyronine Hormones (Lipid-Soluble):} Estrogen, progesterone, testosterone, cortisol, thyroxine ($T_3/T_4$). Cross cell membrane freely; bind to \textbf{intracellular / nuclear receptors} $\implies$ hormone-receptor complex interacts with genome/chromosome to regulate gene expression and protein synthesis.
\end{itemize}
\end{ncertbox}

\end{document}
